\documentclass{../templateReading/cdreading}

\newcommand{\authorinfo}{THE FORUM CENTER - NGUYỄN HOÀNG HUY}
\newcommand{\documenttitle}{TÀI LIỆU CHUYÊN ĐỀ READING}
\newcommand{\collectiontitle}{Level 3 — W7 (Feb 9-15)}

\begin{document}

\thispagestyle{firstpage}
\makeworkshopheader{\authorinfo}{\documenttitle}{\collectiontitle}
\pagestyle{otherpages}

\cdsection{The Grand Banks}

Paragraph A

The Grand Banks is a large area of submerged highlands southeast of Newfoundland and east of the Laurentian Channel on the North American continental shelf. Covering 93,200 square kilometres, the Grand Banks are relatively shallow, ranging from 25 to 100 meters in depth. It is in this area that the cold Labrador Current mixes with the warm waters of the Gulf Stream. The mixing of these waters and the shape of the ocean bottom lifts nutrients to the surface and these conditions created one of the richest fishing grounds in the world. Extensive marine life flourishes in the Grand Banks, whose range extends beyond the Canadian 200-mile exclusive economic zone (EEZ) and into international waters. This has made it an important part of both the Canadian and the high seas fisheries, with fishermen risking their lives in the extremely inhospitable environment consisting of rogue waves, fog, icebergs, sea ice, hurricanes, winter storms and earthquakes.

Paragraph B

While the area’s ‘official’ discovery is credited to John Cabot in 1497, English and Portuguese vessels are known to have first sought out these waters prior to that, based upon reports they received from earlier Viking voyages to Newfoundland. Several navigators, including Basque fishermen, are known to have fished these waters in the fifteenth century. Some texts from that era refer to a land called Bacalao, ‘the land of the codfish’, which is possibly Newfoundland. However, it was not until John Cabot noted the waters’ abundance of sea life that the existence of these fishing grounds became widely known in Europe. Soon, fishermen and merchants from France, Spain, Portugal and England developed seasonal inshore fisheries producing for European markets. Known as ‘dry’ fishery, cod were split, salted, and dried on shore over the summer before crews returned to Europe. The French pioneered ‘wet’ or ‘green’ fishery on the Grand Banks proper around 1550, heavily salting the cod on board and immediately returning home.

Paragraph C

The Grand Banks were possibly the world’s most important international fishing area in the nineteenth and twentieth centuries. Technological advances in fishing, such as sonar and large factory ships, including the massive factory freezer trawlers introduced in the 1950’s, led to overfishing and a serious decline in the fish stocks. Based upon the many foreign policy agreements Newfoundland had entered into prior to its admittance into the Canadian Confederation, foreign fleets, some from as far away as Russia, came to the Grand Banks in force, catching unprecedented quantities of fish.

Paragraph D

Between 1973 and 1982, the United Nations and its member states negotiated the Third Convention of the Law of the Sea, one component of which was the concept of nations being allowed to declare an EEZ. Many nations worldwide-declared 200-nautical mile EEZ’s, including Canada and the United States. On the whole, the EEZ was very well received by fishermen in eastern Canada, because it meant they could fish unhindered out to the limit without fear of competing with the foreign fleets. During the late 1970’s and early 1980s, Canada’s domestic offshore fleet grew as fishermen and fish processing companies rushed to take advantage. It was during this time that it was noticed that the foreign fleets now pushed out to areas of the Grand Banks off Newfoundland outside the Canadian EEZ. By the late 1980’s, dwindling catches of Atlantic cod were being reported throughout Newfoundland and eastern Canada, and the federal government and citizens of coastal regions in the area began to face the reality that the domestic and foreign overfishing had taken its toll. The Canadian government was finally forced to take drastic action in 1992, when a total moratorium was declared indefinitely for the northern cod.

Paragraph E

Over the last ten years, it has been noted that cod appear to be returning to the Grand Banks in small numbers. The reasons for this fragile recovery are still unknown. Perhaps, the damage done by trawlers is not permanent and the marine fauna and ecosystems can rebuild themselves if given a prolonged period of time without any commercial activity. Either way, the early stage recovery of the Grand Banks is encouraging news, but caution is needed, as, after nearly twenty years of severe limitations, cod stocks are still only at approximately ten per cent of 1960’s levels. It is hoped that in another ten to twenty years, stocks may be close to a full recovery, although this would require political pressure to maintain strict limitations on commercial fishing. If cod do come back to the Grand Banks in meaningful numbers, it is to be hoped that the Canadians will not make the same mistakes again.

Paragraph F

Further riches have now been found in the Grand Banks. Petroleum reserves have been discovered and a number of oil fields are under development in the region. The vast Hibernia oil field was discovered in 1979, and, following several years of aborted start-up attempts, the Hibernia megaproject began construction of the production platform and gravity base structures in the early 1990’s. Production commenced on November 17, 1997, with initial production rates in excess of 50,000 barrels of crude oil per day from a single well. Hibernia has proven to be the most prolific oil well in Canada. However, earthquake and iceberg activity in the Grand Banks pose a potential ecological disaster that could devastate the fishing grounds that are only now starting to recover.

\cdsection{Questions 14-20}

The text on the previous pages has 6 paragraphs A – F.

Which paragraph contains the following information?

\cdnoteinline{Write your answers in boxes 14-20 on your answer sheet.}

\textbf{14} Many countries could legally fish Newfoundland waters because of treaties Newfoundland had made before becoming part of Canada.

\textbf{15} The establishment of the EEZ did not stop over-fishing in the Grand Banks.

\textbf{16} Natural disasters could cause oil to destroy what is left of the Grand Banks ecosystem.

\textbf{17} The original amount of fish in the Grand Banks was due to different temperature waters mixing.

\textbf{18} East Canadian fishermen were generally happy with the establishment of the Canadian EEZ.

\textbf{19} Grand Banks’ cod stocks are still 90 per cent lower than what they were in the 1960’s.

\textbf{20} The French were the first to prepare the cod on board their ships before going back to France.

\cdsection{Questions 21-23}

\cdnoteinline{Choose the correct letter A, B, C or D.}

\cdnoteinline{Write the correct letter in boxes 21-23 on your answer sheet.}

\textbf{21} The first English fishermen to come to the Grand Banks to fish

A   were told about the fishery by Basque fishermen.

B   were sent word about the fishery from the first American colonists.

C   acted on information from previous Viking expeditions.

D   discovered the fishery themselves while exploring.

\textbf{22} John Cabot’s reports of the Grand Banks

A   led to the establishment of the Canadian EEZ.

B   meant the fishery was well known in Europe.

C   led to fighting between rival fishing fleets.

D   were not immediately publicised, so that English fishermen could benefit.

\textbf{23} The establishment of the Canadian EEZ

A   did not stop foreign fishermen from fishing the Grand Banks.

B   was not ratified by the United Nations.

C   temporarily stopped the over-fishing of cod in the Grand Banks.

D   meant Canadian fishermen were excluded from fishing the Grand Banks.

\cdsection{Questions 24-26}

Do the following statements agree with the information given in the text?

In boxes 24-26 on your answer sheet write:

TRUE               if the statement agrees with the information

FALSE              if the statement contradicts the information

NOT GIVEN    if there is no information on this

\textbf{24} Even now, cod stocks have shown no signs of recovery in the Grand Banks.

\textbf{25} Initial efforts to extract oil from the Grand Banks’ Hibernia oil field were unsuccessful.

\textbf{26} Oil exploration companies have to follow strict safety controls imposed by the Canadian government.

\cdblue{\textbf{Hướng dẫn giải}}

Câu 14

Nội dung: Many countries could legally fish Newfoundland waters because of treaties Newfoundland had made before becoming part of Canada.
\begin{enumerate}
\item Từ khóa tìm trong bài
\begin{itemize}
\item treaties
\item Newfoundland
\item before becoming part of Canada
\end{itemize}
\item Từ khóa mấu chốt
\begin{itemize}
\item Đoạn C: "...foreign policy agreements Newfoundland had entered into prior to its admittance into the Canadian Confederation..."
\end{itemize}
\item Phân tích \& Vị trí
\begin{itemize}
\item "Treaties" (hiệp ước) tương đương với "foreign policy agreements" (thỏa thuận chính sách đối ngoại).
\item "Before becoming part of Canada" tương đương với "prior to its admittance into the Canadian Confederation". Đoạn văn giải thích lý do các hạm đội nước ngoài được phép đánh bắt cá.
\end{itemize}
\item Đáp án: C
\end{enumerate}

Câu 15

Nội dung: The establishment of the EEZ did not stop over-fishing in the Grand Banks.
\begin{enumerate}
\item Từ khóa tìm trong bài
\begin{itemize}
\item establishment of the EEZ
\item did not stop
\item over-fishing
\end{itemize}
\item Từ khóa mấu chốt
\begin{itemize}
\item Đoạn D: "...foreign fleets now pushed out to areas... outside the Canadian EEZ... domestic and foreign overfishing had taken its toll."
\end{itemize}
\item Phân tích \& Vị trí
\begin{itemize}
\item Đoạn văn đề cập đến việc thiết lập EEZ (vùng đặc quyền kinh tế), nhưng sau đó chỉ ra rằng các hạm đội nước ngoài vẫn đánh bắt ngay bên ngoài vùng này và tình trạng "overfishing" (đánh bắt quá mức) vẫn gây thiệt hại ("taken its toll"). Điều này có nghĩa là việc thiết lập EEZ không ngăn chặn được vấn đề.
\end{itemize}
\item Đáp án: D
\end{enumerate}

Câu 16

Nội dung: Natural disasters could cause oil to destroy what is left of the Grand Banks ecosystem.
\begin{enumerate}
\item Từ khóa tìm trong bài
\begin{itemize}
\item Natural disasters
\item oil
\item destroy ecosystem
\end{itemize}
\item Từ khóa mấu chốt
\begin{itemize}
\item Đoạn F: "...earthquake and iceberg activity... pose a potential ecological disaster that could devastate the fishing grounds..."
\end{itemize}
\item Phân tích \& Vị trí
\begin{itemize}
\item "Earthquake and iceberg activity" (động đất và băng trôi) là các thảm họa tự nhiên ("Natural disasters").
\item "Devastate" (tàn phá) đồng nghĩa với "destroy". "Ecological disaster" tương ứng với việc hủy hoại hệ sinh thái ("ecosystem").
\end{itemize}
\item Đáp án: F
\end{enumerate}

Câu 17

Nội dung: The original amount of fish in the Grand Banks was due to different temperature waters mixing.
\begin{enumerate}
\item Từ khóa tìm trong bài
\begin{itemize}
\item original amount of fish
\item different temperature waters mixing
\end{itemize}
\item Từ khóa mấu chốt
\begin{itemize}
\item Đoạn A: "...cold Labrador Current mixes with the warm waters... created one of the richest fishing grounds in the world."
\end{itemize}
\item Phân tích \& Vị trí
\begin{itemize}
\item "Richest fishing grounds" (ngư trường giàu có nhất) ám chỉ lượng cá ban đầu ("original amount of fish").
\item "Cold... mixes with... warm waters" chính là sự pha trộn của các dòng nước có nhiệt độ khác nhau ("different temperature waters mixing").
\end{itemize}
\item Đáp án: A
\end{enumerate}

Câu 18

Nội dung: East Canadian fishermen were generally happy with the establishment of the Canadian EEZ.
\begin{enumerate}
\item Từ khóa tìm trong bài
\begin{itemize}
\item East Canadian fishermen
\item happy
\item establishment of the Canadian EEZ
\end{itemize}
\item Từ khóa mấu chốt
\begin{itemize}
\item Đoạn D: "On the whole, the EEZ was very well received by fishermen in eastern Canada..."
\end{itemize}
\item Phân tích \& Vị trí
\begin{itemize}
\item Cụm từ "very well received" (được đón nhận rất tốt) mang ý nghĩa tương đương với việc ngư dân cảm thấy vui vẻ/hài lòng ("happy").
\end{itemize}
\item Đáp án: D
\end{enumerate}

Câu 19

Nội dung: Grand Banks’ cod stocks are still 90 per cent lower than what they were in the 1960’s.
\begin{enumerate}
\item Từ khóa tìm trong bài
\begin{itemize}
\item cod stocks
\item 90 per cent lower
\item 1960’s
\end{itemize}
\item Từ khóa mấu chốt
\begin{itemize}
\item Đoạn E: "...cod stocks are still only at approximately ten per cent of 1960’s levels."
\end{itemize}
\item Phân tích \& Vị trí
\begin{itemize}
\item Bài đọc nói trữ lượng cá tuyết chỉ còn khoảng 10\% so với mức của những năm 1960 ("ten per cent of 1960’s levels"). Điều này đồng nghĩa với việc nó thấp hơn 90\% ("90 per cent lower").
\end{itemize}
\item Đáp án: E
\end{enumerate}

Câu 20

Nội dung: The French were the first to prepare the cod on board their ships before going back to France.
\begin{enumerate}
\item Từ khóa tìm trong bài
\begin{itemize}
\item French
\item first
\item prepare cod on board
\item before going back
\end{itemize}
\item Từ khóa mấu chốt
\begin{itemize}
\item Đoạn B: "The French pioneered ‘wet’ or ‘green’ fishery... heavily salting the cod on board and immediately returning home."
\end{itemize}
\item Phân tích \& Vị trí
\begin{itemize}
\item "Pioneered" (tiên phong) tương ứng với "first".
\item "Salting the cod on board" (muối cá trên tàu) là hành động sơ chế ("prepare cod on board").
\item "Immediately returning home" tương ứng với "going back to France".
\end{itemize}
\item Đáp án: B
\end{enumerate}

Câu 21

Nội dung: The first English fishermen to come to the Grand Banks to fish...
\begin{enumerate}
\item Từ khóa tìm trong bài
\begin{itemize}
\item first English fishermen
\item come to the Grand Banks
\end{itemize}
\item Từ khóa mấu chốt
\begin{itemize}
\item Đoạn B: "...English and Portuguese vessels are known to have first sought out these waters... based upon reports they received from earlier Viking voyages..."
\end{itemize}
\item Phân tích \& Vị trí
\begin{itemize}
\item Ngư dân Anh đến vùng này dựa trên các báo cáo từ các chuyến đi trước đó của người Viking ("reports... from earlier Viking voyages"). Điều này khớp với lựa chọn C ("acted on information from previous Viking expeditions").
\end{itemize}
\item Đáp án: C
\end{enumerate}

Câu 22

Nội dung: John Cabot’s reports of the Grand Banks...
\begin{enumerate}
\item Từ khóa tìm trong bài
\begin{itemize}
\item John Cabot’s reports
\end{itemize}
\item Từ khóa mấu chốt
\begin{itemize}
\item Đoạn B: "...John Cabot noted the waters’ abundance of sea life that the existence of these fishing grounds became widely known in Europe."
\end{itemize}
\item Phân tích \& Vị trí
\begin{itemize}
\item Nhờ ghi chép của John Cabot mà sự tồn tại của ngư trường này trở nên "widely known in Europe" (được biết đến rộng rãi ở Châu Âu). Điều này khớp với lựa chọn B.
\end{itemize}
\item Đáp án: B
\end{enumerate}

Câu 23

Nội dung: The establishment of the Canadian EEZ...
\begin{enumerate}
\item Từ khóa tìm trong bài
\begin{itemize}
\item establishment
\item Canadian EEZ
\end{itemize}
\item Từ khóa mấu chốt
\begin{itemize}
\item Đoạn D: "...foreign fleets now pushed out to areas of the Grand Banks off Newfoundland outside the Canadian EEZ."
\end{itemize}
\item Phân tích \& Vị trí
\begin{itemize}
\item Việc thiết lập EEZ khiến các hạm đội nước ngoài bị đẩy ra ngoài vùng đặc quyền, nhưng họ vẫn đánh bắt ở khu vực Grand Banks nằm ngoài ranh giới đó. Tức là EEZ không ngăn được ngư dân nước ngoài đánh bắt tại Grand Banks (chỉ là đánh bắt ở phần ngoài EEZ). Điều này khớp với lựa chọn A ("did not stop foreign fishermen from fishing the Grand Banks").
\end{itemize}
\item Đáp án: A
\end{enumerate}

Câu 24

Nội dung: Even now, cod stocks have shown no signs of recovery in the Grand Banks.
\begin{enumerate}
\item Từ khóa tìm trong bài
\begin{itemize}
\item cod stocks
\item no signs of recovery
\end{itemize}
\item Từ khóa mấu chốt
\begin{itemize}
\item Đoạn E: "...cod appear to be returning to the Grand Banks in small numbers... early stage recovery of the Grand Banks is encouraging news..."
\end{itemize}
\item Phân tích \& Vị trí
\begin{itemize}
\item Bài đọc khẳng định cá đang quay trở lại ("returning") và sự phục hồi giai đoạn đầu là tin tức đáng khích lệ ("encouraging news"). Thông tin này mâu thuẫn hoàn toàn với mệnh đề "no signs of recovery" (không có dấu hiệu phục hồi).
\end{itemize}
\item Đáp án: FALSE
\end{enumerate}

Câu 25

Nội dung: Initial efforts to extract oil from the Grand Banks’ Hibernia oil field were unsuccessful.
\begin{enumerate}
\item Từ khóa tìm trong bài
\begin{itemize}
\item Initial efforts
\item extract oil
\item Hibernia oil field
\item unsuccessful
\end{itemize}
\item Từ khóa mấu chốt
\begin{itemize}
\item Đoạn F: "...the Hibernia oil field was discovered in 1979, and, following several years of aborted start-up attempts..."
\end{itemize}
\item Phân tích \& Vị trí
\begin{itemize}
\item "Aborted start-up attempts" (các nỗ lực khởi động bị hủy bỏ) đồng nghĩa với việc những nỗ lực ban đầu ("initial efforts") là không thành công ("unsuccessful").
\end{itemize}
\item Đáp án: TRUE
\end{enumerate}

Câu 26

Nội dung: Oil exploration companies have to follow strict safety controls imposed by the Canadian government.
\begin{enumerate}
\item Từ khóa tìm trong bài
\begin{itemize}
\item Oil exploration companies
\item strict safety controls
\item Canadian government
\end{itemize}
\item Từ khóa mấu chốt
\begin{itemize}
\item Đoạn F (toàn bộ đoạn nói về dầu khí).
\end{itemize}
\item Phân tích \& Vị trí
\begin{itemize}
\item Đoạn F đề cập đến việc phát hiện dầu, sản lượng và rủi ro thiên tai, nhưng không có thông tin nào nói về các quy định an toàn nghiêm ngặt ("strict safety controls") do chính phủ Canada áp đặt lên các công ty này.
\end{itemize}
\item Đáp án: NOT GIVEN
\end{enumerate}

\end{document}
